\documentclass[xcolor=pdftex,dvipsnames,table]{beamer}


\mode<presentation> {
  \usetheme{Warsaw}
  \setbeamercovered{transparent}
%  \useoutertheme{infolines}
%  \setbeamertemplate{headline}[default]
%  \setbeamertemplate{footline}[infolines theme]{}
}

\usepackage[english]{babel}

\usepackage[latin1]{inputenc}

\usepackage{times}
\usepackage[T1]{fontenc}

\usepackage{ulem}

\hypersetup{%
    pdftitle={Stealth Distributed Hash Table: A Robust and Flexible Super-Peered DHT},
    pdfauthor={Andrew Brampton, Andrew MacQuire, Idris A. Rai, Nicholas J. P. Race and Laurent Mathy},
    pdfkeywords={Distributed Hash Tables, Peer-to-Peer, Stealth DHT},
    bookmarksnumbered,
    pdfstartview={FitV},
    linkcolor={black},%: Color for normal internal links.
    anchorcolor={black},%: Color for anchor text.
    citecolor={black},%: Color for bibliographical citations in text.
    filecolor={black},%: Color for URLs which open local files.
    menucolor={black},%: Color for Acrobat menu items.
    pagecolor={black},%: Color for links to other pages
    urlcolor={black},%: Color for linked URLs.
}%

\title[Stealth Distributed Hash Table]
{\textbf{Stealth Distributed Hash Table}\\A Robust and Flexible Super-Peered DHT}

\institute[Lancaster University, UK] {
  \textbf{Andrew Brampton}, Andrew MacQuire, Idris A. Rai, \\Nicholas J. P. Race and Laurent Mathy\\
  {\it \{brampton,macquire,rai,race,laurent\}@comp.lancs.ac.uk}\\
  Computing Department\\
  Lancaster University, UK
}

\date[7th December 2006]
{Thursday 7th December \\ \textbf{CoNEXT 2006}}

\subject{Stealth Distributed Hash Table}

\logo{\includegraphics[height=0.85cm]{uni-logo-win}}

\setbeamertemplate{navigation symbols}{}

%\renewcommand{\textbf}{\alert}

\begin{document}

\begin{frame}
  \titlepage
\end{frame}

\begin{frame}
  \frametitle{Outline}
  \tableofcontents
\end{frame}

\section{Introduction}
%\subsection{Distributed Hash Table Overview}
\subsection{Distributed Hash}
%\begin{frame}
%  \frametitle{Distributed Hash Table Overview}
%  \begin{columns}
%
%  \column{6cm}
%  \begin{itemize}
%
%  \item{\textbf{DHTs are:}}
%     \begin{itemize}
%     \item{Scalable}
%     \item{Decentralised}
%     \item{Self-organising}
%     \item{Fault-tolerant}
%     \item{Overlay networks}
%     \end{itemize}
%
%  ~\\
%  \item{Examples}
%    \begin{itemize}
%        \item{\textbf{Pastry}}
%        \item{Chord}
%        \item{CAN}
%        \item{Tapestry}
%    \end{itemize}
%
%
%  \hfill\\
%
%%  \item<2-5>{\textbf{Capable of:}}
%%     \begin{itemize}
%%     \item{{\it Efficient} request routing}
%%     \item{{\it Deterministic} object location}
%%     \item{{\it Effective} load balancing}
%%     \end{itemize}
%
%  \end{itemize}
%
%  \column{6cm}
%
%  \begin{center}
%    \begin{overlayarea}{\textwidth}{0.8\textheight}
%        \includegraphics<1>[width=6cm]{diagrams/DHT}
%%         \includegraphics<3>[width=6cm]{diagrams/DHT-1}
%%         \includegraphics<4>[width=6cm]{diagrams/DHT-2}
%%         \includegraphics<5>[width=6cm]{diagrams/DHT-3
%    \end{overlayarea}
%  \end{center}
%
%\end{columns}
%\end{frame}

\begin{frame}
  \frametitle{Distributed Hash Table Overview}
  \begin{columns}

  \column{6cm}
    \textbf{Pastry}
     \begin{itemize}
     \item{Circular Identifier Space (typically $0-2^{128}$)}
     \item{Each node has ID\\Each key has ID}
     \item{Nodes are responsible for keys closest to them (on the ring)}
     \item{Routing done iteratively closer and closer to a ID}
     \end{itemize}
  \column{6cm}
  \begin{center}
    \begin{overlayarea}{\textwidth}{0.8\textheight}
        \includegraphics<1>[width=6cm]{diagrams/DHT}
    \end{overlayarea}
  \end{center}
  \end{columns}

\end{frame}

\begin{frame}
  \frametitle{Distributed Hash Table Problems}
  \begin{columns}

  \column{6cm}

  \begin{itemize}
    \frametitle{}
    \item<1-2>{\textbf{Assumed homogeneity}}~\\
        {Varying bandwidth, processing or storage abilities}
    \item<3>{\textbf{Churn overhead}}~\\
        {Joining causes updates, Leaving causes staleness}
    \item<4>{\textbf{Free-riding}}~\\
        {Users don't want to participate, or stay for a very short time}
    \item<5>{\textbf{Security issues}}~\\
        {Sniffing, Man in the middle, Poisoning, Denial of Service}
  \end{itemize}

  \column{6cm}

  \begin{center}
    \begin{overlayarea}{\textwidth}{0.8\textheight}
        \includegraphics<1>[width=6cm]{diagrams/DHT}
        \includegraphics<2>[width=6cm]{diagrams/Homogeneity}
        \includegraphics<3-4>[width=6cm]{diagrams/Churn}
        %AB: Maybe some kind of graphic for 4
        \includegraphics<5>[width=5.5cm]{diagrams/Security}
    \end{overlayarea}
  \end{center}
\end{columns}
\end{frame}


\subsection{Related Work}
\begin{frame}
    \frametitle{Related Work}

    \begin{columns}

    \column{6cm}

    \begin{center}
        \begin{overlayarea}{\textwidth}{0.8\textheight}
        \includegraphics<1>[width=6cm]{diagrams/StaticBinding}

        \includegraphics<2>[width=6cm]{diagrams/MultipleOverlays}
        \end{overlayarea}
    \end{center}

    \column{5cm}
    Typical approaches to tackling outlined problems:\\~\\
    \begin{itemize}
    \item<1>{\textbf{Static binding}}
        \begin{itemize}
        \item{Single point of failure per client}
        %Peer-to-Peer Internet Telephony using SIP - Singh et al
        \item{\it e.g. Singh et al.}
        \end{itemize}
    \item<2>{\textbf{Multiple overlays}}
        \begin{itemize}
        \item{Large associated overhead}
        %A Super-Peer Based Lookup in Structured Peer-to-Peer Systems - Zhu et al
        %Structured Superpeers: Leveraging Heterogeneity to Provide Constant-Time Lookup - Mizrak et al
        \item{\it e.g. Mizrak et al. and Zhu et al.}
        \end{itemize}
    \end{itemize}

\end{columns}
\end{frame}

\section{Stealth DHT}
\begin{frame}
    \frametitle{Stealth DHT}

    \begin{itemize}
    \item{Provides support for \textbf{heterogeneous} nodes on a DHT}
    \end{itemize}

    \begin{columns}

    \column{6cm}

    \begin{itemize}
    \item<1>{\textbf{Service Nodes} (Super-peers)}
        \begin{itemize}
        \item{Assumed to be \textbf{highly provisioned}}
        \item{Responsible for \textbf{forwarding and storing} data}
        \item{Just a normal DHT node}
        \end{itemize}
    \item<2>{\textbf{Stealth Nodes}}
        \begin{itemize}
        \item{Potentially \textbf{underpowered or unreliable} ({\it e.g.} mobile)}
        \item{\textbf{No} forwarding or storing responsibilities}
        \item{\textbf{Invisible} to all nodes, \textbf{including} Service nodes}
        \end{itemize}
    \end{itemize}

    \column{6cm}

    \begin{overlayarea}{\textwidth}{0.8\textheight}
    \includegraphics<1>[width=6cm]{Diagrams/ServiceHighlight}
    \includegraphics<2>[width=6cm]{Diagrams/StealthHighlight}
    \end{overlayarea}

    \end{columns}

\end{frame}

\subsection{Join Procedures}
\begin{frame}
    \frametitle{
    \only<1-10>{Service Node / Pastry Join Procedure}
    \only<11-12>{Stealth Node Join Procedure}
    }

    \begin{columns}

    \column{5cm}

    \only<1-9>{

    \textbf{Phase A:} Gathering state
        \begin{itemize}
        \item<1>{Node \it{X} wants to join}
        \item<2>{\it{X} contacts bootstrap \it{A}}
        \item<3>{\it{A} responds with state}
        \item<4>{\it{A} forwards towards \it{Y}}
        \item<5>{More state returned}
        \item<6>{Continues towards \it{Y}}
        \item<7>{More state returned}
        \item<8>{Reaches closest node \it{Y}}
        \item<9>{\it{Y} informs \it{X} of finish}
        \end{itemize}
    }
    \only<10>{

    \textbf{Phase B:} Announcement
        \begin{itemize}
        \item{X announces presence}
        \item{Others update tables}
        \end{itemize}
        \hfill\\
        \hfill\\
    X is now a member of the DHT and can be used for \textbf{forwarding
    messages}, \textbf{storing data} etc.
    }

    \only<11>{
    \textbf{Stealth DHT Phase A}
        \begin{itemize}
        \item{A stealth node X gathers state in the same way as before}
        \end{itemize}
    }

    \only<12>{
    \textbf{Stealth nodes \it{do not} announce their presence}
        \begin{itemize}
        \item{This ensures stealth nodes \textbf{cannot} be used for forwarding messages or storing data}
        \end{itemize}
    }

\column{6cm}
  \begin{center}
    \begin{overlayarea}{\textwidth}{0.8\textheight}
        \includegraphics<1>[width=6cm]{diagrams/PhaseA-1}
        \includegraphics<2>[width=6cm]{diagrams/PhaseA-2}
        \includegraphics<3>[width=6cm]{diagrams/PhaseA-3}
        \includegraphics<4>[width=6cm]{diagrams/PhaseA-4}
        \includegraphics<5>[width=6cm]{diagrams/PhaseA-5}
        \includegraphics<6>[width=6cm]{diagrams/PhaseA-6}
        \includegraphics<7>[width=6cm]{diagrams/PhaseA-7}
        \includegraphics<8>[width=6cm]{diagrams/PhaseA-8}
        \includegraphics<9>[width=6cm]{diagrams/PhaseA-9}
        \includegraphics<10>[width=6cm]{diagrams/PhaseB}
        \includegraphics<11>[width=6cm]{diagrams/PhaseA-Stealth}
        \includegraphics<12>[width=6cm]{diagrams/PhaseB-No}
    \end{overlayarea}
  \end{center}
  \end{columns}
\end{frame}

\subsection{Routing Information}
\begin{frame}
    \frametitle{Routing Information}

    \begin{columns}

    \column{5.5cm}
    \begin{overlayarea}{\textwidth}{1\textheight}
    Service Node (0x0F23...)

    \begin{itemize}
        \item{Routing table}
    \end{itemize}

    \rowcolors[\hline]{1}{blue!20}{blue!10}
    \begin{tabular}{l!{\vrule}c c c c c c}
    & 0 & 1 & 2 & 3 & ... & F\\
    \hline
    0 & - & x & x & x & ... & x\\
    1 & x & x & x & x & ... & -\\
    2 & x & - & x & x & ... & x\\
    . & . & . & . & . & ... & .\\
    \end{tabular}

    ~\\
    \begin{itemize}
        \item{Leafset}
    \end{itemize}

    \rowcolors[\hline]{1}{blue!20}{blue!10}
    \begin{tabular}{c c c c c c c}
    \hline
    x & x & x & ... & x & x & x\\
    \hline
    \end{tabular}

    \end{overlayarea}

    \column{5.5cm}

    \begin{center}
        \begin{overlayarea}{\textwidth}{\textheight}
            \includegraphics[width=6cm]{diagrams/ServiceHighlight}
        \end{overlayarea}
    \end{center}

    \end{columns}
\end{frame}

\begin{frame}
    \frametitle{Routing Information}

    \begin{columns}

    \column{5.5cm}
    \begin{overlayarea}{\textwidth}{1\textheight}
    Stealth Node (0x3284...)

    \begin{itemize}
        \item{Routing table}
        \begin{itemize}
            \item{Single row}
            \item{Extra cell}
            \item{\footnotesize 97\% of time first row is used}
            \item{\footnotesize Better than random first hops}
        \end{itemize}
    \end{itemize}

    \rowcolors[\hline]{1}{blue!20}{blue!10}
    \begin{tabular}{l!{\vrule}c c c c c c}
    \hline
      & 0 & 1 & 2 & 3 & ... & F\\
    0 & x & x & x & x & ... & x\\
    \end{tabular}

    ~\\
    \begin{itemize}
        \item{No Leafset}
    \end{itemize}

    \end{overlayarea}

    \column{5.5cm}

    \begin{center}
        \begin{overlayarea}{\textwidth}{\textheight}
            \includegraphics[width=6cm]{diagrams/StealthHighlight}
        \end{overlayarea}
    \end{center}

    \end{columns}

\end{frame}

\subsection{Stealth Node State}
\begin{frame}
    \frametitle{Stealth Node State}

    \begin{itemize}
    \item{Stealth node routing tables become \textbf{stale}...}
    \end{itemize}

    \begin{columns}

    \column{6cm}

    \begin{overlayarea}{\textwidth}{0.75\textheight}
        \includegraphics<1>[width=6cm]{diagrams/Polling}
        \includegraphics<2>[width=6cm]{diagrams/Piggybacking}
        \includegraphics<3>[width=6cm]{diagrams/Rejoin}
    \end{overlayarea}

    \column{6cm}

    \begin{itemize}
        \item<1>{\textbf{Polling}}
            \begin{itemize}
            \item{Periodically requesting state}
            \end{itemize}
            ~\\
        \item<2>{\textbf{Piggybacking}}
            \begin{itemize}
            \item{Receiving state with replies}
            \end{itemize}
            ~\\
        \item<3>{\textbf{Rejoining}}
            \begin{itemize}
            \item{Occasional full rejoin}
            \end{itemize}
    \end{itemize}

    \end{columns}
\end{frame}

%\begin{frame}
%    \frametitle{Recovery Method Comparison}
%    \includegraphics[width=10.75cm]{diagrams/Fig2a}
%\end{frame}
%
%\begin{frame}
%    \frametitle{Recovery Method Comparison}
%    \includegraphics[width=10.75cm]{diagrams/Fig2b}
%\end{frame}

\subsection{Stealth Node Addressability}
\begin{frame}
    \begin{itemize}
    \item{Stealth nodes \textbf{cannot be contacted} by default...}
    \end{itemize}

    \frametitle{Stealth Node Addressability}
    \begin{itemize}
        \item<1>{\textbf{Registration}}
        \begin{itemize}
        \item{Stealth nodes \textbf{register} with their \textbf{closest} service node}
        \item{Service nodes forward reply messages to their registered stealth nodes}
        \end{itemize}

    \hfill\\

    \item<2>{\textbf{Encapsulation}}
        \begin{itemize}
        \item{Stealth nodes route a message towards the destination.}
        \item{The first Service nodes \textbf{encapsulate} the Stealth node's network details in the outgoing message and act as a NAT}
        \end{itemize}
    \end{itemize}

\end{frame}

\begin{frame}
    \frametitle{Addressability Method Comparison}
    \begin{center}
%    Dependent nodes in an Application Level Multicast example
    \includegraphics[width=10cm]{diagrams/Addressability}
    \end{center}
\end{frame}

\begin{frame}
    \frametitle{Benefits}
    \begin{itemize}
    \item{Nodes on the same DHT are \textbf{separated based on capabilities}}
        \begin{itemize}
        \item{\textbf{More powerful} nodes are given \textbf{more responsibility}}
        \end{itemize}
    ~\\
    \item{Churning stealth nodes generate \textbf{little to no} overhead}
        \begin{itemize}
        \item{Stealth node joins are \textbf{lightweight} i.e. No Announcements}
        \item{Service node have \textbf{no} Stealth node \textbf{state}}
        \end{itemize}
    ~\\
    \item{\textbf{Improved security}}
        \begin{itemize}
        \item{If service nodes are \textbf{appropriately authenticated} then \textbf{only authorised} nodes store and forward data}
        \item{Stealth nodes \textbf{cannot} commit \textbf{sniff}ing or corruption attacks}
        \end{itemize}
%    \item{Support for commercial applications}
%        \begin{itemize}
%        \item{Service nodes provided by a commercial entity}
%        \end{itemize}
    \end{itemize}
\end{frame}

\section{Evaluation}

\subsection{Methodology}
\begin{frame}
    \frametitle{Evaluation Methodology}
    \begin{itemize}
%    \item{\textbf{Simulator}}
%        \begin{itemize}
%        \item{Our own Java-based discrete-event packet-level platform}
%        \end{itemize}
%        ~\\
%    \item{\textbf{Implementation}}
%        \begin{itemize}
%        \item{Our own C++ based system, running over PlanetLab}
%        \end{itemize}
%        ~\\
    \item{\textbf{Methodology}}
        \begin{itemize}
        \item{Simulations \& PlanetLab implementation}
        \item{Pastry as reference (and results shown relative to Pastry)}
        \end{itemize}
        ~\\
    \item{\textbf{Typical Workload}}
        \begin{itemize}
        \item{Large number of keys placed into the network}
        \item{Key popularity according to a Zipf distribution}
        \item{With and without churn}
        \end{itemize}
    \end{itemize}
\end{frame}

% results section

\subsection{Results}
%\begin{frame}
%    \frametitle{Validation}
%    \includegraphics[width=10.75cm]{diagrams/Validation}
%\end{frame}

\begin{frame}
    \frametitle{Join Overhead: Messages (Simulation)}
    \includegraphics[width=10.75cm]{diagrams/Fig4}
\end{frame}

\begin{frame}
    \frametitle{Join Overhead: Messages (PlanetLab)}
    \includegraphics[width=10.75cm]{diagrams/Fig13}
\end{frame}

%\begin{frame}
%    \frametitle{Join Overhead: Latency (Simulation)}
%    \includegraphics[width=10.75cm]{diagrams/Fig5}
%\end{frame}

\begin{frame}
    \frametitle{Lookup: Latency with and without churn (Simulation)}
    \includegraphics[width=10.75cm]{diagrams/Fig6}
\end{frame}

\begin{frame}
    \frametitle{Lookup: Misses under churn (Simulation)}
    \includegraphics[width=10.75cm]{diagrams/Fig7}
\end{frame}

%\begin{frame}
%    \frametitle{Lookup: Misses under decreasing churn (Simulation)}
%    \includegraphics[width=10.75cm]{diagrams/Fig11}
%\end{frame}

\begin{frame}
    \frametitle{Load balancing: No Churn (Simulation)}
    \includegraphics[width=10.75cm]{diagrams/Fig8a}
\end{frame}

\begin{frame}
    \frametitle{Load balancing: Churn (Simulation)}
    \includegraphics[width=10.75cm]{diagrams/Fig8b}
\end{frame}

\begin{frame}
    \frametitle{Single row routing vs Random selection}
    \includegraphics[width=10.75cm]{diagrams/RandomNodeStressLog}
\end{frame}

\subsection{Conclusion}
\begin{frame}
    \frametitle{Conclusion}

    % Better conclusion/summary

    \begin{itemize}
        \item{Pros}
        \begin{itemize}
            \item{Reduces effects of clients churning}
            \item{Increase performance due to smaller routing population}
            \item{Improves security}
        \end{itemize}
        ~\\
        \item{Cons}
        \begin{itemize}
            \item{Service nodes churn has greater impact}
        \end{itemize}
        ~\\
        \item{Applications}
        \begin{itemize}
            \item{Mobile environments}
            \item{Content Distribution Networks (CDNs)}
        \end{itemize}
    \end{itemize}

%    \begin{itemize}
%    \item{Stealth DHTs:}
%        \begin{itemize}
%        \item{\textbf{Partition} the network into two groups}
%        \item{\textbf{Increase performance} in most areas}
%        \item{\textbf{Return control} to the service provider}
%        \end{itemize}
%    \end{itemize}

%    \begin{itemize}
%    {\footnotesize
%    \item{Extended Work}
%        \begin{itemize}
%        {\footnotesize
%        \item{{\it Performance Analysis of Stealth DHT with Mobile Nodes} - MP2P'06}
%        \item{{\it Authentication in Stealth Distributed Hash Tables} - Euromicro'06}
%        \item{{\it Performance Modelling of Peer-to-Peer Routing} - Submitted to HotP2P'07}
%        }
%        \end{itemize}
%    }
%    \end{itemize}
\end{frame}

\begin{frame}
  %\frametitle{Thank you for listening}
  \begin{center}
    \textbf{Thank you for listening\\
    Any questions?\\}
    \titlepage
  \end{center}
\end{frame}

\end{document}
